\section{One explicit language and its hard slices}
\label{sec:construction}

The language will evaluate a polynomial at a linear hash of its data bits.
The hash and the polynomial coefficients are themselves input bits.  This
lets us choose them probabilistically when proving a lower bound, while
keeping the membership algorithm fixed.  We first define that algorithm on
every input, then prove that suitable settings of its non-data bits give
small correlation with every bounded-width CNF on most random restrictions.
The input-index construction in~\cite[Section~5]{PSZ00} supplies an earlier
instance of this evaluator-and-hard-subfunction method.  Its lower-bound
model has bottom fan-in two; the correlation property proved below is the
one needed for arbitrary fixed CNF width.

\begin{definition}[The language]\label{def:language}
For an input of length $n$, put $d=\floor{n/5}$ and reject if $d<64$.
Otherwise define
\[
 r=\ceil{d^{2/3}},\qquad
 t=\max\{j\in\N\cup\{0\}: j\text{ is even and }j^6\le d\}.
\]
Reject if fewer than $2d+(t+2)r-1$ input bits are available.
Otherwise parse the input into the following consecutive blocks:
\begin{enumerate}
\item $d$ data bits $x=(x_1,\ldots,x_d)$;
\item $d+r-1$ hash bits $u_1,\ldots,u_{d+r-1}$;
\item $r$ bits $p_0,\ldots,p_{r-1}$, encoding the monic polynomial
\[
 P(Z)=Z^r+\sum_{i=0}^{r-1}p_iZ^i\in\F_2[Z];
\]
\item $t$ blocks of $r$ bits, encoding elements $\beta_0,\ldots,\beta_{t-1}$ of
the quotient ring $R_P=\F_2[Z]/(P)$, in the basis represented by
$1,Z,\ldots,Z^{r-1}$.
\end{enumerate}
Unused bits are ignored.  Form
\begin{equation}\label{eq:toeplitz-hash}
 h_i(x)=\sum_{s=1}^d u_{i+s-1}x_s\pmod2
 \quad(1\le i\le r),\qquad
 h(x)=\sum_{i=1}^r h_i(x)Z^{i-1}\in R_P.
\end{equation}
Let $\lambda:R_P\to\F_2$ take the degree-zero coefficient of the unique
representative of degree less than $r$.  The input belongs to $L$ exactly
when
\begin{equation}\label{eq:language-evaluation}
 \lambda\left(\sum_{j=0}^{t-1}\beta_jh(x)^j\right)=0.
\end{equation}
\end{definition}

The matrix in \eqref{eq:toeplitz-hash} is Hankel (constant on
anti-diagonals), or Toeplitz after reversing the data coordinates.
Fully random Toeplitz matrices give a standard universal linear hash
family~\cite[pp.~136--137]{Krawczyk94}.  We use only the uniform image
of each nonzero difference, proved below, not strong pairwise
independence of hash values.  Once the hash is injective on a restricted
cube, polynomial interpolation supplies limited-independent values,
as in~\cite[Section~6, proof of Proposition~6.3]{ABI86}.

The rejection rule fixes the language at the small lengths once and for
all.  For $d\ge64$, the blocks always fit: they use
\begin{equation}\label{eq:construction-block-length}
 d+(d+r-1)+r+tr=2d+(t+2)r-1
\end{equation}
bits, and $r\le2d^{2/3}$ and $t+2\le2d^{1/6}$ imply
$(t+2)r\le4d^{5/6}\le2d$.  Thus
\eqref{eq:construction-block-length} is at most $4d-1\le n$.
We also have $t\ge2$ and
\begin{equation}\label{eq:construction-growth}
 d^{1/6}-2<t\le d^{1/6},\qquad
 d^{2/3}\le r<d^{2/3}+1.
\end{equation}

\begin{proposition}\label{prop:polytime}
The language $L$ is decided on every input of length $n$ by one
deterministic Turing machine in time $C(n+1)^a$, for some fixed positive
integers $C,a$.
\end{proposition}

\begin{proof}
After the small-length check, the integer parameters can be computed
without real arithmetic: $r$ is the least positive integer with
$r^3\ge d^2$, and $t$ can be found by testing even integers until their
sixth powers exceed $d$.  Searching integers up to $n$ and using ordinary
integer arithmetic takes polynomial time.  Reading and checking the
block boundaries does also.

Since $P$ is monic, division by $P$ gives a unique remainder of degree
less than $r$, whether or not $P$ is irreducible.  Thus every allowed
polynomial block gives well-defined ring operations.  Direct multiplication
of two coefficient arrays and reduction by $P$ use $O(r^2)$ bit operations:
form a polynomial of degree at most $2r-2$, then successively cancel its
highest coefficients by shifted copies of $P$.  Addition uses $O(r)$ bit
operations.  Formula~\eqref{eq:toeplitz-hash} uses $O(rd)$ bit operations,
and Horner evaluation of \eqref{eq:language-evaluation} uses $O(t)$ ring
multiplications and additions.  The evaluation therefore uses
$O(rd+tr^2)$ bit-array operations and polynomial storage.

These are fixed loops on finite bit arrays, with counters of polynomial
size.  They can be implemented by a fixed Turing machine using direct tape
scans, which incur only polynomial overhead per array operation.  Together
with the initial integer computations this yields a polynomial time bound;
increasing its coefficient and exponent to positive integers gives the
stated $C,a$ and covers the finitely many small lengths as well.  In
particular, the machine does not test irreducibility or search for any
choice of parameter bits.  Neither its instructions nor its running-time
constants depend on the proof parameters used below.
\end{proof}

For $d\ge64$ and fixed non-data bits, write $g:\bits^d\to\{-1,1\}$ for
the sign of the resulting data function, with $+1$ on accepted inputs.
It equals $(-1)^{\lambda(\sum_j\beta_jh(x)^j)}$.
Although the evaluator allows every monic modulus, we will select an
irreducible polynomial of degree $r$ for the hardness analysis.  Such a
polynomial exists for every $r\ge1$; an elementary proof is given in
Lemma~\ref{lem:field-existence} in Appendix~\ref{sec:field-existence}.
The next proposition is the property of these slices that the circuit
argument will use.

\begin{proposition}\label{prop:good-slices}
For each fixed integer $b\ge1$ and real $B>0$, all sufficiently large
lengths $n$ admit fixed settings of the non-data bits such that, with
\[
 p=\frac{B}{\sqrt d},\qquad m_0=pd=B\sqrt d,
\]
the resulting sign $g$ satisfies
\begin{equation}\label{eq:good-slices}
 \Pr_{T,\rho}\left[
 |T|\notin[m_0/2,2m_0]
 \ \text{or}\ 
 \corr_b(g_{T,\rho})>2^{-|T|/8}
 \right]
 \le4\cdot2^{-m_0/16}.
\end{equation}
Here $T$ includes each of the $d$ data coordinates independently with
probability $p$, and $\rho$ is an independent uniform assignment outside
$T$.  In particular, the correlation bound on a good restriction holds
simultaneously for all members of $\mathcal C_b(T)$.
\end{proposition}

\begin{proof}
Fix $b,B$ and take $d$ large enough that $0<p<1$ and $t\ge2$.
Choose a monic irreducible polynomial $P$ of degree $r$, and identify
$R_P$ with its field $F$ of size $2^r$.  Initially
choose all the hash bits $u$ independently and uniformly, and choose
$\beta_0,\ldots,\beta_{t-1}$ independently and uniformly in $F$.  We estimate
the failure probability jointly over these choices and the restriction;
the final step will fix one choice of these input bits.

\paragraph{Injectivity on a restricted cube.}
For any nonzero $w\in\bits^d$, the vector $h(w)$ is uniform in
$\F_2^r$ over the hash bits.  Indeed, let $s$ be the largest coordinate
with $w_s=1$.  The $i$th row in \eqref{eq:toeplitz-hash} contains
$u_{i+s-1}$ with coefficient one, and every other bit it uses has a
smaller index.  In particular no earlier row uses $u_{i+s-1}$.  The $r$
rows therefore have independent successive pivots, so the linear map
$u\mapsto h(w)$ has rank $r$ and takes a uniform input to a uniform
output.  Equivalently, after choosing the other bits, these pivot bits
can be set successively in exactly one way to give any prescribed output.
Thus $\Pr_u[h(w)=0]=2^{-r}$.

Fix a restriction with $m=|T|$.  Its affine data cube has exactly
$2^m-1$ nonzero differences, namely the nonzero vectors supported on $T$.
Since $h$ is linear, two cube points collide exactly when their nonzero
difference hashes to zero.  A union bound gives
\begin{equation}\label{eq:hash-collision}
 \Pr_u[h\text{ is not injective on the restricted cube}]
 \le(2^m-1)2^{-r}\le2^{m-r}.
\end{equation}
The fixed assignment $\rho$ only translates the hash image and has no
effect on this estimate.

\paragraph{Limited independence from interpolation.}
Condition on a hash that is injective on this cube.  At any $q\le t$
distinct cube points, let $a_1,\ldots,a_q\in F$ be their distinct hashes.
For arbitrary desired values $v_1,\ldots,v_q\in F$, the polynomial
\[
 Q(X)=\sum_{i=1}^q v_i
       \prod_{\substack{1\le j\le q\\j\ne i}}
          \frac{X-a_j}{a_i-a_j}
\]
has degree at most $q-1<t$ and satisfies $Q(a_i)=v_i$.
The denominators are nonzero field elements.  Consequently the linear map
from $(\beta_0,\ldots,\beta_{t-1})\in F^t$ to the $q$ evaluation values is
surjective.  Each fiber is a translate of its kernel and hence has the
same size, so uniform coefficients give independent uniform values in
$F^q$.  The coordinate map $\lambda$ is a nonzero $\F_2$-linear map,
with $\lambda(1)=1$.  Its two fibers have equal size, for translation by
$1$ interchanges them.  Applying $(-1)^\lambda$ coordinatewise therefore
gives unbiased, $t$-wise independent signs on the restricted cube.

\paragraph{One bound for every sparse test.}
Apply Lemma~\ref{lem:sparse-partition} with $\eta=1/8$ and write
$M=M(b,1/8)$ for its positive integer constant.  On a cube of dimension
$m\ge1$, the number of possible normalized clauses of width at most $b$,
including the false clause of width zero, is
\[
 C_m=\sum_{j=0}^{\min(b,m)}2^j\binom mj\le(2m+1)^b.
\]
For the inequality, list a clause's literals in a fixed order and pad
to length $b$ with a symbol outside the $2m$ literals; this is an
injective encoding.  A CNF with at most $Mm$ clauses can in turn be
encoded by a list of length $Mm$ over these $C_m$ clauses and a padding
symbol.  In particular the number of such CNF indicators is at most
\begin{equation}\label{eq:sparse-test-count}
 (C_m+1)^{Mm}
 \le 2^{Mm[1+b\log_2(2m+1)]}.
\end{equation}
This list includes the constant true test (no clauses) and the constant
false test (one empty clause).

Fix one such test $J$, and write
$Y=\sum_{z\in\bits^T}g_{T,\rho}(z)J(z)$.  Expectation in the next
calculation is over the coefficients $\beta_j$, conditional on the injective
hash.  Expand $Y^t$ as a sum over ordered $t$-tuples of cube points.
If some point occurs exactly once, its sign is independent of all the
other distinct signs in the term and has mean zero, so that term's
expectation vanishes.  Otherwise at most $t/2$ distinct points occur.
All such tuples are covered by listing $t/2$ cube points, with repetitions
allowed to pad the list, and choosing a position in that list for each
of the $t$ factors.  There are at most
$2^{mt/2}(t/2)^t$ such descriptions.  Each term has absolute expectation
at most one, and hence
\begin{equation}\label{eq:even-moment}
 \E Y^t\le t^t2^{mt/2}.
\end{equation}
Because $t$ is even, $Y^t=|Y|^t$.  Markov's inequality, obtained by
bounding $|Y|^t$ below by a threshold on its exceedance event, gives
\begin{equation}\label{eq:single-test-tail}
 \Pr\bigl[|Y|>2^{3m/4}\bigr]
 \le t^t2^{-mt/4}.
\end{equation}
By \eqref{eq:sparse-test-count}, the probability that any sparse test
has $|\E_z g_{T,\rho}(z)J(z)|>2^{-m/4}$ is at most
\begin{equation}\label{eq:all-sparse-tail}
 2^{Mm[1+b\log_2(2m+1)]+t\log_2t-mt/4}
 \le2^{-mt/8}
\end{equation}
whenever
\begin{equation}\label{eq:sparse-union-condition}
 \frac{M[1+b\log_2(2m+1)]}{t}
 +\frac{\log_2t}{m}\le\frac18.
\end{equation}
This condition holds uniformly for $m_0/2\le m\le2m_0$ once $d$ is
sufficiently large.  Indeed, with $b,B,M$ fixed,
\eqref{eq:construction-growth} bounds the two summands respectively by
$O_{b,B,M}(\log d/d^{1/6})$ and $O_B(\log d/\sqrt d)$, both tending
to zero.

If no sparse test fails, any $J\in\mathcal C_b(T)$ is, by
Lemma~\ref{lem:sparse-partition}, a sum of at most $2^{m/8}$ indicators
from the sparse list.  The triangle inequality thus gives
\begin{equation}\label{eq:all-width-tests}
 \abs{\E_z g_{T,\rho}(z)J(z)}
 \le2^{m/8}\,2^{-m/4}=2^{-m/8}.
\end{equation}
The sparse event already controls every indicator in the list, so this
conclusion holds for every width-$b$ CNF simultaneously.  No further
probability estimate depends on the choice of $J$.

\paragraph{Restriction tails and a fixed setting.}
The variable $m=|T|$ is binomial with mean $m_0$.  Independence and
$1+v\le e^v$ give
\[
 \E 2^{-m}=(1-p/2)^d\le e^{-m_0/2},\qquad
 \E 2^m=(1+p)^d\le e^{m_0}.
\]
Markov's inequality in these two estimates yields
\begin{align}
 \Pr[m<m_0/2]
 &\le2^{-((\log_2e-1)/2)m_0},\label{eq:live-lower-tail}\\
 \Pr[m>2m_0]
 &\le2^{-(2-\log_2e)m_0}.\label{eq:live-upper-tail}
\end{align}
Both displayed exponents have coefficient greater than $1/16$.
For example, tangent and secant bounds for $1/x$ on $[1,2]$ give
$2/3<\ln2<3/4$, and therefore $4/3<\log_2e<3/2$.
The total probability of leaving the prescribed range is at most
$2\cdot2^{-m_0/16}$.

For every restriction in that range, \eqref{eq:hash-collision} and
\eqref{eq:all-sparse-tail} bound the conditional failure probability over
the parameter bits by $2^{m-r}+2^{-mt/8}$.  Uniformly in this range,
\[
 2^{m-r}\le2^{-m_0/16},\qquad
 2^{-mt/8}\le2^{-m_0/16}
\]
for all sufficiently large $d$: the first inequality follows from
$r/m_0\to\infty$, by taking $r\ge(2+1/16)m_0$, and the second from
$m\ge m_0/2$ and $t\ge1$.  Adding these two errors and the two
restriction tails gives joint failure probability at most
$4\cdot2^{-m_0/16}$.

For each fixed choice of $(u,\beta_0,\ldots,\beta_{t-1})$, let its failure
probability be the probability over $(T,\rho)$ of the event in
\eqref{eq:good-slices}.  The average of these finitely many numbers is
the joint failure probability just bounded.  At least one setting
therefore has failure probability at most $4\cdot2^{-m_0/16}$.
Fix that setting, the already selected $P$, and arbitrary unused bits.
This proves the proposition.
\end{proof}

The choices in Proposition~\ref{prop:good-slices} may depend on $b,B$
and the input length.  The quantifiers require only that $b,B$ be fixed
before the length tends to infinity.  These choices select restrictions
of the single language in Definition~\ref{def:language}; computing
membership never requires finding them.
